\documentclass[UTF8]{ctexart}
\usepackage{amsmath}
\usepackage{graphicx}
\usepackage{booktabs}
\usepackage{hyperref}

\title{卷积神经网络在图像识别中的应用综述}
\author{何思颖 \\ 人工智能一班 126230501}
\date{\today}

\begin{document}

\maketitle

\begin{abstract}
本文综述了卷积神经网络的基本原理与发展历程，介绍了卷积层、池化层的核心思想，
并通过 MNIST 手写数字识别任务验证其有效性。实验结果表明，CNN 在该任务上
可达 99\% 以上准确率。
\end{abstract}

\section{引言}
深度学习近年来发展迅速\cite{lecun2015}，卷积神经网络（CNN）因其参数共享和局部
连接的特性，在图像识别领域取得了突破性进展。

\subsection{研究背景}
传统图像识别依赖手工设计特征，泛化能力有限。CNN 通过端到端学习自动提取特征，
显著提升了识别精度。

\section{方法}
\subsection{卷积操作}
卷积运算可表示为公式~\eqref{eq:conv}：
\begin{equation}
    y_i = \sum_j w_j x_{i+j} + b
    \label{eq:conv}
\end{equation}

\subsection{网络结构}
本文采用经典 LeNet-5 结构，包含两个卷积层与两个全连接层。

\section{实验}
\subsection{数据集}
实验使用 MNIST 手写数字数据集，部分样本如图~\ref{fig:sample} 所示。

\begin{figure}[htbp]
    \centering
    \includegraphics[width=0.35\textwidth]{sample.jpg}
    \caption{MNIST 样本示例}
    \label{fig:sample}
\end{figure}

\subsection{实验结果}
不同方法的准确率对比如表~\ref{tab:res} 所示。

\begin{table}[htbp]
    \centering
    \caption{不同方法在 MNIST 上的准确率对比}
    \begin{tabular}{lcc}
        \toprule
        方法 & 准确率 & 参数量 \\
        \midrule
        SVM & 97.8\% & 少 \\
        MLP & 98.5\% & 中 \\
        CNN & 99.2\% & 多 \\
        \bottomrule
    \end{tabular}
    \label{tab:res}
\end{table}

\section{结论}
本文介绍了 CNN 的核心原理并通过实验验证其有效性。未来将探索其在更复杂
视觉任务中的应用。

\begin{thebibliography}{9}
    \bibitem{lecun2015} LeCun Y, Bengio Y, Hinton G. Deep learning[J]. Nature, 2015, 521: 436-444.
    \bibitem{he2016} He K, Zhang X, Ren S, et al. Deep residual learning for image recognition[C]//CVPR. 2016: 770-778.
\end{thebibliography}

\end{document}
